\SourceLeaf{185}\clearpage
\begin{center}[ASTRALGEOMETRIE.]\end{center}

\begin{center}\rule{3cm}{0.4pt}\end{center}

\begin{center}[\textsc{Gerling} \emph{an} \textsc{Gauss}. \emph{Marburg, 23. Juli 1818.}]\end{center}

\{.\ .\ .\ .\ .\ .\ Mir steht im künftigen Semester eine Arbeit bevor, für welche ich
schon jetzt so frei bin Sie um Rath zu bitten, und diess zu thun,
wahrscheinlich öfter so frei sein werde. --- \textsc{Fleckeisen} hat mich
nemlich gebeten, die Besorgung einer neuen Auflage von \textsc{Lorenz}’
reiner Mathematik zu übernehmen; und ich kann diess um so lieber thun, da
ich hier nach diesem Buch selbst lese; und es nun schon 6 Jahre immer beim
Unterricht gebraucht habe.

.\ .\ .\ .\ .\ .\ .\ In der Geometrie habe ich weniger Anstösse, möchte aber besonders
gern Ihre Meinung wissen, wie es mit der Parallelentheorie wohl am besten zu
halten ist. Was \textsc{Lorenz} darüber hat, ist theils falsch, theils
ungründlich. --- Dass die Euklidische Manier vorzutragen sei, dabei aber der
Mangel derselben einzugestehen, halte ich für Recht; das quomodo aber ist
mir nicht klar. Ich habe gedacht, es sei am besten, den Satz: Eine gerade
Linie kann durch einen Punkt nur eine Parallele haben, als Axiom
hinzustellen, und in einer Anmerkung zu sagen, dass man einen Beweis für
diesen Satz noch nicht habe finden können, und deshalb ihn so lange bis
einer gefunden, oder die Unwahrheit des Satzes bewiesen sei, als Axiom
annehmen müsse, wie im Grunde schon Euklides gethan. --- Haben Sie doch die
Güte mir auch hierüber Ihr Urtheil zu eröffnen. .\ .\ .\ .\ .\ .\}

\begin{center}\rule{3cm}{0.4pt}\end{center}

